% Транскрипция рабочих заметок с фотографий 184041 и 184042.
% Зачёркнутые фрагменты опущены; сокращения раскрыты.
% Номерные метки связывают определения с опорными точками.
% Вспомогательные определения набраны с отступом слева.
% Сборка из корня репозитория: working/1a-graphs/_src/render.sh
\documentclass[12pt,a4paper]{article}
\usepackage[margin=20mm]{geometry}
\usepackage{fontspec}
\setmainfont{Times New Roman}
\usepackage{polyglossia}
\setdefaultlanguage{russian}
\usepackage{amsmath,amssymb}
\usepackage{xcolor}
\usepackage{fancyhdr}
\pagestyle{fancy}
\fancyhf{}
\lhead{ДЗ 1a}
\rhead{Опорные точки}
\cfoot{\thepage}
\setlength{\headheight}{15pt}
\setlength{\parindent}{0pt}
\setlength{\parskip}{6pt}
\newcommand{\Aut}{\operatorname{Aut}}
\newcommand{\term}[1]{\par\textbf{#1.}\enspace}
\definecolor{anchorink}{HTML}{24516B}
\definecolor{anchorfill}{HTML}{EAF2F7}
\newcommand{\anchorlabel}[1]{\begingroup
  \setlength{\fboxsep}{4pt}\setlength{\fboxrule}{0.4pt}%
  \fcolorbox{anchorink}{anchorfill}{\color{anchorink}\small\bfseries №\,#1}%
  \endgroup}
\newcommand{\anchorterm}[2]{\par\addvspace{8pt}%
  \noindent\anchorlabel{#1}\quad\textbf{#2.}\par\nobreak\noindent}
\newenvironment{auxiliary}{\begin{list}{}{%
  \setlength{\leftmargin}{14mm}\setlength{\rightmargin}{0pt}%
  \setlength{\topsep}{3pt}\setlength{\parsep}{5pt}%
  \setlength{\partopsep}{0pt}}\item\relax}{\end{list}}

\begin{document}
\begin{center}
{\Large\bfseries Опорные точки}
\end{center}

\begin{enumerate}
\item Граф $\longrightarrow$ неориентированный граф.
\item Изоморфность графов (и неизоморфность графов)
  $\longrightarrow$ изоморфизм.
\item Валентность вершины $\longrightarrow$ отношение инцидентности.
\item Полнота списка (что больше нет?) $\longrightarrow$ да!
\item $\Gamma$ --- «нарисованные» графы (?) $\longrightarrow$ да!
\item $|\Aut\Gamma|$ --- ?
  \begin{itemize}
  \item $\Aut\Gamma\longrightarrow\{\ldots\}\longrightarrow$ автоморфизм;
  \item $|\ldots|$ --- мощность множества.
  \end{itemize}
\item Как через это доказать полноту?
  \begin{itemize}
  \item масса графа/системы;
  \item полная система;
  \item масса системы;
  \item монотонность.
  \end{itemize}
\end{enumerate}

\medskip\hrule\medskip
\section*{Раскрытие понятий}
\anchorterm{1}{Граф}
Граф --- это тройка из множества вершин, множества рёбер и отображения,
ставящего в соответствие каждому ребру пару вершин, упорядоченную или
неупорядоченную.

\begin{auxiliary}
\term{Ориентированный граф}
Это граф, у каждого ребра которого есть начало и конец.

\term{Неориентированный граф}
Это граф, у каждого ребра которого есть два равноправных конца.
\end{auxiliary}

\anchorterm{2}{Изоморфные графы}
Два графа, между которыми существует хотя бы один изоморфизм.

\begin{auxiliary}
\term{Изоморфизм графов}
Это пара биекций: вершин на вершины и рёбер на рёбра, согласованные
с отображениями инцидентности обоих графов.
\end{auxiliary}

\anchorterm{3}{Валентность вершины}
Это количество рёбер, которые присоединены к вершине.

\begin{auxiliary}
\term{Отношение инцидентности}
Соответствие между вершиной и ребром, концом которого она является.
Приписка: «входит», «входит/выходит» (орграф), «присоединено».
\end{auxiliary}

\newpage
\section*{Автоморфизмы, масса и полнота}
\anchorterm{6}{Автоморфизм графа}
Это изоморфизм графа самого на себя.

\anchorterm{7}{Масса графа}
Число, обратное количеству его автоморфизмов.

\anchorterm{7}{Масса системы}
Это масса множества всех попарно неизоморфных графов, валентности которых
равны компонентам вектора:
\[
  E=\sum_{\Gamma}\frac{1}{|\Aut\Gamma|}.
\]
Это масса минимальной полной системы графов, валентности которых равны
компонентам вектора.

\anchorterm{7}{Монотонность массы по системе графов}
\[
\text{Масса системы}\leq\text{масса полной системы},
\]
\[
\text{равенство}\quad\Longleftrightarrow\quad\text{система полная}.
\]
Так как масса монотонна относительно системы попарно неизоморфных графов.

\medskip\hrule\medskip
\section*{Формула массы системы}
\[
\boxed{E=\sum_{\Gamma}\frac{1}{|\Aut\Gamma|}=\frac{PQ}{(2N)!}}
\]
$(n_1,n_2,\ldots,n_k)$ --- вектор валентностей.
\[
N=\frac12\sum_i n_i,
\qquad
P=\frac{(2N)!}{N!\,2^N},
\]
\[
Q=\frac{(2N)!}
  {(n_1!)^{m_1}m_1!\,(n_2!)^{m_2}m_2!\cdots}.
\]
$m_i$ --- частота $n_i$ в векторе валентностей.
\end{document}
